This section develops the pricing framework. \Cref{sec:model-setup} states the one property every
model must respect. \Cref{sec:model-m0,sec:model-m1,sec:model-m2} build the single-contract model
in three layers: a Gaussian ``Bachelier digital'' with a deterministic information clock (M0), a
stochastic-volatility extension with a general backbone (M1), and jumps (M2).
\Cref{sec:model-ladder} treats strips of contracts on one numeric outcome, and
\Cref{sec:model-hazard} treats contracts that pay if an event occurs before a deadline. Proofs are
in \cref{app:derivations}.

\subsection{Setup}\label{sec:model-setup}

Fix a filtered probability space $(\Omega,\F,(\F_t)_{t\le T},\Prob)$. A binary contract pays one
unit at $T$ if an event $A\in\F_T$ occurs. Its fair price is the conditional probability
\begin{equation}\label{eq:fair}
  \pr_t=\Prob(A\mid\F_t),\qquad 0\le t\le T .
\end{equation}
We abstract from interest (collateral is posted in full and settlement horizons are short), from
fees and from risk premia. Whether market prices equal $\pr_t$ is an empirical question we return
to in \cref{sec:discussion}. By construction $\pr_t$ is a martingale bounded in $[0,1]$ with
$\pr_T\in\{0,1\}$. That alone fixes the total variance.

\begin{lemma}[Variance identity]\label{lem:variance}
For any binary contract, $\E\big[\langle \pr\rangle_T\big]=\E\big[(\pr_T-\pr_0)^2\big]=\pr_0(1-\pr_0)$,
where $\langle\pr\rangle$ is the quadratic variation. For prices observed at
$0=t_0<t_1<\dots<t_n=T$, $\E\sum_i(\pr_{t_{i+1}}-\pr_{t_i})^2=\pr_0(1-\pr_0)$.
\end{lemma}

A model of prediction-market prices is therefore a model of \emph{when} and \emph{how} a fixed
budget of variance $\pr_0(1-\pr_0)$ is released. It specifies the timing (a clock), the dependence
on the price level (a backbone), randomness in the release rate (stochastic volatility) and its
lumpiness (jumps). These are the four dimensions the empirical work measures.

\subsection{M0: the Bachelier digital and the information clock}\label{sec:model-m0}

\begin{definition}[M0]\label{def:m0}
A latent score $Y$ satisfies $\dd Y_t=s(t)\,\dd W_t$ with a deterministic information flow
$s(\cdot)$, and $A=\{Y_T>0\}$. Write
\[
  \Vrem_t=\int_t^T s(u)^2\,\dd u
\]
for the remaining information variance, and $\irate_t=s(t)^2/\Vrem_t$ for the
\emph{information rate}, the fraction of the remaining uncertainty resolved per unit time.
\end{definition}

The construction goes back to \citet{stern1994}'s Brownian model of sports scores. Its implication
for the volatility of a probability price was developed by \citet{archakipeirotis2009}, and
\citet{taleb2018} applies it to election forecasts.

\begin{proposition}[Pricing, dynamics and transition under M0]\label{prop:m0}
Under M0, with $\Phi$ and $\varphi$ the standard normal distribution and density:
\begin{enumerate}
  \item \emph{Price.} $\pr_t=\Phi\big(Y_t/\sqrt{\Vrem_t}\big)$. As a function $\pr(t,y)$ it solves the
  backward heat equation with a digital payoff,
  \begin{equation}\label{eq:pde}
    \partial_t \pr+\tfrac12 s(t)^2\,\partial_{yy}\pr=0,\qquad \pr(T,y)=\ind{y>0}.
  \end{equation}
  This is the Black--Scholes equation with zero rate and normal (Bachelier) dynamics for a
  cash-or-nothing payoff.
  \item \emph{Dynamics.} In probit units $\Zt_t=\Phi^{-1}(\pr_t)=Y_t/\sqrt{\Vrem_t}$,
  \begin{equation}\label{eq:zdyn}
    \dd\Zt_t=\sqrt{\irate_t}\,\dd W_t+\tfrac12\irate_t\Zt_t\,\dd t,\qquad
    \dd\pr_t=\varphi(\Zt_t)\sqrt{\irate_t}\,\dd W_t .
  \end{equation}
  \item \emph{Transition.} For $h>0$ let $R=\Vrem_t/\Vrem_{t+h}\ge1$. Then
  \begin{equation}\label{eq:transition}
    \Zt_{t+h}\mid\F_t\sim\N\!\big(\Zt_t\sqrt R,\;R-1\big),\qquad
    \Prob(\pr_{t+h}\le x\mid\pr_t)=\Phi\!\left(\frac{\Phi^{-1}(x)-\Phi^{-1}(\pr_t)\sqrt R}{\sqrt{R-1}}\right).
  \end{equation}
\end{enumerate}
\end{proposition}

Three implications organise the empirical work.
\begin{itemize}
  \item \emph{Clock.} The price's instantaneous volatility is $\varphi(\Phi^{-1}(\pr))\sqrt{\irate_t}$:
  a level factor times a clock factor. With a uniform flow ($s$ constant) the clock is
  $\irate_t=1/\tauexp$, with $\tauexp=T-t$, so volatility at a given price rises like
  $\tauexp^{-1/2}$ into expiry.
  \item \emph{Test statistic.} The standardised innovation
  \begin{equation}\label{eq:u}
    \ustd=\frac{\Zt_{t+h}-\Zt_t\sqrt R}{\sqrt{R-1}}
  \end{equation}
  is standard normal under M0. Its variance tests the clock and its kurtosis tests the Gaussian
  shape.
  \item \emph{Drift.} In probit units the price drifts away from zero at rate $\tfrac12\irate\Zt$. The
  martingale requires this drift to offset the convexity of $\Phi$.
\end{itemize}

The analogues of the option Greeks follow directly. Delta is $\partial\pr/\partial Y=\varphi(\Zt)/\sqrt{\Vrem}$.
The ``vega'' with respect to remaining variance is $\partial\pr/\partial\Vrem=-\tfrac12\Zt\varphi(\Zt)/\Vrem$:
more uncertainty pulls the price toward $\tfrac12$. Theta at a fixed score is
$\partial\pr/\partial t=\tfrac12\Zt\varphi(\Zt)\irate$: with no news, favourites drift up and
longshots down. Theta is exactly offset in expectation by the diffusion term of \eqref{eq:pde}.

\subsection{M1: a bounded SABR model}\label{sec:model-m1}

M0 ties volatility to the level through $\varphi(\Phi^{-1}(\pr))$ and to time through a
deterministic clock. M1 relaxes both, in the spirit of SABR \citep{hagan2002}.

\begin{definition}[M1]\label{def:m1}
\begin{equation}\label{eq:m1}
  \dd\pr_t=\alpha_t\,c(\tauexp_t)\,g_\beta(\pr_t)\,\dd W_t,\qquad
  g_\beta(\pr)=[\pr(1-\pr)]^\beta,\qquad
  \dd\alpha_t=\nu\,\alpha_t\,\dd B_t,\qquad
  \dd\langle W,B\rangle_t=\rho\,\sgn(\Zt_t)\,\dd t .
\end{equation}
\end{definition}

The parameters play the following roles.
\begin{itemize}
  \item \emph{Backbone $\beta$.} It nests the leading bounded diffusions:
  \begin{itemize}
    \item $\beta=1$ is the logit-normal model, $\dd\,\mathrm{logit}\,\pr=\sigma\,\dd W+\dots$, the
    analogue of lognormal dynamics;
    \item $\beta=\tfrac12$ is the Wright--Fisher or Jacobi diffusion \citep{karlintaylor1981}, the
    analogue of square-root dynamics;
    \item the M0 backbone $\varphi(\Phi^{-1}(\pr))$ behaves like $\pr\sqrt{2\ln(1/\pr)}$ near zero,
    between the two.
  \end{itemize}
  Pure Bachelier dynamics for $\pr$ ($\beta=0$) are inadmissible, because they do not keep the
  price in $[0,1]$.
  \item \emph{Clock $c(\cdot)$.} The deterministic information schedule.
  \item \emph{Vol of vol $\nu$.} How strongly the information rate itself fluctuates.
  \item \emph{Correlation $\rho$.} The factor $\sgn(\Zt)$ makes the specification invariant to
  relabelling YES and NO. With $\rho<0$, volatility rises when the price moves toward
  $\tfrac12$; with $\rho>0$, it rises as the price moves toward resolution.
\end{itemize}

\Cref{lem:variance} holds for every specification of M1. Stochastic volatility redistributes the
variance budget across time and paths; it does not enlarge it.

\subsection{M2: jumps}\label{sec:model-m2}

Information often arrives discretely: scores, rulings, releases, the start of a speech. M2 adds
compensated jumps to the latent score,
\begin{equation}\label{eq:m2}
  \dd Y_t=s(t)\,\dd W_t+J_t\,\dd N_t,
\end{equation}
where $N$ is a counting process with intensity $\kappa(t)$ and $J$ an independent jump size. With
Gaussian jumps $J\sim\N(\mu_J,\sigma_J^2)$, a constant intensity and $n$ jumps remaining with
Poisson probability $\pi_n$, the price is a \citet{merton1976}-type mixture,
\begin{equation}\label{eq:m2price}
  \pr_t=\sum_{n\ge0}\pi_n\,\Phi\!\left(\frac{Y_t+n\mu_J}{\sqrt{\Vrem_t+n\sigma_J^2}}\right).
\end{equation}
Between jumps, the drift that keeps $\pr$ a martingale is the compensator of the expected jump,
not the diffusive convexity term. The M0 link between local variance and outward probit drift
(\cref{prop:m0}, item 2) therefore breaks when a share of the variance arrives in jumps. We use
this in \cref{sec:res-hf}.

\subsection{Strips of digitals on one outcome: implied distributions and smiles}\label{sec:model-ladder}

Many contracts are listed as ladders on one numeric outcome $X_T$, such as a temperature, a
points total or a score. Each leg pays $\ind{X_T>K}$ (or $\ind{K_1<X_T\le K_2}$ for brackets). The
survival function $S_t(K)=\Prob(X_T>K\mid\F_t)$ is then the price of a strip of digital options,
and it identifies the conditional distribution of $X_T$. This is the digital counterpart of
\citet{breedenlitzenberger1978}.

\begin{definition}[Implied median, implied volatility, moneyness]\label{def:smile}
The implied median is $\Fmed_t=S_t^{-1}(\tfrac12)$. The Bachelier-digital implied volatility of
strike $K$ solves
\[
  S_t(K)=\Phi\big((\Fmed_t-K)/\sigma_K\big),
\]
and moneyness is $m=(K-\Fmed_t)/\satm$. The smile is the map $m\mapsto\sigma_K/\satm$, summarised by
\[
  \sigma_K/\satm-1=a+s\,m+c\,m^2 ,
\]
with skew $s$ and curvature $c$.
\end{definition}

\begin{proposition}[Shape of the digital smile]\label{prop:smile}
\begin{enumerate}
  \item If $X_T\mid\F_t$ is Gaussian, the smile is flat: $\sigma_K\equiv\satm$.
  \item If $X_T\mid\F_t$ is symmetric with tails heavier than the Gaussian of the same
  interquartile range, then $\sigma_K>\satm$ for every strike beyond the quartiles: the smile turns
  up in the wings. For Student-$t$ outcomes, $\sigma_K/\satm$ increases in $|m|$ (verified
  numerically; \cref{app:derivations}).
  \item Skewness tilts the smile: right skew makes $\sigma_K$ larger above $\Fmed$ than below it.
\end{enumerate}
\end{proposition}

Two further concepts carry over from equity options. The first is smile dynamics: a smile is
\emph{sticky-strike} if $\sigma_K$ at a fixed $K$ does not respond to moves in $\Fmed$, and
\emph{sticky-moneyness} if $\sigma_K$ depends on $K$ only through $m$ \citep{derman1999}. The
second is a variance-swap analogue. If $X_T$ is Gaussian given $\F_t$, each leg is an M0 contract
with $Y=X-K$ and $\Vrem_t=\satm(t)^2$. The ladder's at-the-money term structure is then an
\emph{implied information clock}, and the implied median is a martingale with
\begin{equation}\label{eq:varswap}
  \E\big[(\Fmed_{t+h}-\Fmed_t)^2\big]=\E\big[\satm(t)^2-\satm(t+h)^2\big].
\end{equation}
Comparing the two sides is the prediction-market analogue of comparing realized and implied
variance \citep{carrwu2009,bollerslev2009}.

\subsection{Deadline contracts: a reduced-form hazard model}\label{sec:model-hazard}

A large class of contracts pays if an event occurs \emph{by} a date $T$: a product release, a
statement, an elimination, a weather system forming. We model the event time $\vartheta$ with an
intensity $h_t$, as in reduced-form credit models
\citep{jarrowturnbull1995,lando1998,duffiesingleton1999}.

\begin{proposition}[Hazard pricing]\label{prop:hazard}
On $\{\vartheta>t\}$ the price is
\[
  \pr_t=1-\E_t\Big[\exp\Big(-\int_t^T h_u\,\dd u\Big)\Big].
\]
\begin{enumerate}
  \item \emph{Deterministic hazard.} $\pr_t=1-e^{-H(t,T)}$ with $H(t,T)=\int_t^T h_u\,\dd u$. Before
  the event the price decays deterministically,
  \begin{equation}\label{eq:theta}
    \frac{\dd\pr_t}{\dd t}=-h_t\,(1-\pr_t),
  \end{equation}
  and at $\vartheta$ it jumps to $1$. The decay is exactly the compensator of the jump, so $\pr$ is
  a martingale.
  \item \emph{Implied average hazard.} $\hbarimp_t=-\ln(1-\pr_t)/(T-t)$.
  \item \emph{Term structure.} For a strip with deadlines $T_1<T_2$ written on the same event, the
  forward hazard is
  \[
    h(T_1,T_2)=-\frac{\ln\big[(1-\pr_t(T_2))/(1-\pr_t(T_1))\big]}{T_2-T_1}\ \ge 0,
  \]
  so $\pr_t(T_1)\le\pr_t(T_2)$ is a static no-arbitrage condition.
\end{enumerate}
\end{proposition}

A Cox specification, in which $h$ is itself a diffusion (for example the square-root intensity
of \citealp{cox1985cir}), adds news-driven diffusion to the deterministic decay and the resolution jump. We call the
deterministic decay in \eqref{eq:theta} the contract's theta.

\paragraph{Summary.} The three contract types share one structure: a latent state, an information
clock that releases a fixed variance budget, and jumps. They differ in the latent state, which is
a Gaussian score, a numeric outcome, or an event time. \Cref{tab:model-map} maps the objects to
their option-pricing analogues.

\begin{table}[t]
\centering
\begin{threeparttable}
\caption{Model objects and their option-pricing analogues.}\label{tab:model-map}
\small
\begin{tabular}{@{}lll@{}}
\toprule
Prediction-market object & Symbol & Analogue \\
\midrule
Remaining information variance & $\Vrem_t$ & total implied variance $\sigma^2\tauexp$ \\
Information rate (clock) & $\irate_t=s_t^2/\Vrem_t$ & forward variance per unit of total variance \\
Probit price & $\Zt=\Phi^{-1}(\pr)$ & normalised moneyness $d_2$ \\
Backbone & $\beta$ in $[\pr(1-\pr)]^\beta$ & CEV exponent / SABR $\beta$ \\
Vol state, vol of vol, correlation & $\alpha,\nu,\rho$ & SABR $\alpha,\nu,\rho$ \\
Ladder implied median & $\Fmed$ & forward \\
Ladder implied volatility & $\sigma_K$, $\satm$ & Bachelier implied volatility \\
Deadline implied hazard & $\hbarimp$ & credit spread / default intensity \\
Decay with no event & $-h(1-\pr)$ & carry of a short credit position \\
\bottomrule
\end{tabular}
\end{threeparttable}
\end{table}
